\documentclass[12pt]{article}
\usepackage[utf8]{inputenc}
\usepackage[T1]{fontenc}
\usepackage{hyperref}
\title{Pisaflores Tepehua Language Resources for K-12 Curriculum Development}
\author{Compiled by Claude Sonnet 5 (claude-sonnet-5) for Andr\'as Kornai}
\date{}
\begin{document}
\maketitle

\section{Introduction}
Pisaflores Tepehua (ISO 639-3: tpp; Glottocode pisa1237) is a Totonacan language of Mexico with roughly 4,000 speakers. It is endangered. No school programme using it is known. Web search found little material, and nothing written for classroom use. Everything listed below is documentation, scripture or reference material.

\section{Primary Educational Resources}
Search found no primers, textbooks or school materials in Pisaflores Tepehua.

\section{Linguistic and Reference Materials}
\begin{itemize}
\item NSF grant BCS-0132293 (Carolyn MacKay and Frank Trechsel) aimed to produce a grammatical description and a trilingual Tepehua/Spanish/English lexicon. The search results do not show whether either was published: \url{https://grantome.com/grant/NSF/BCS-0132293}
\item MacKay's FAMSI-funded project (2002--2003) built a corpus with a lexicon, about 1000 diagnostic phrases, and narratives and conversations. A planned sketch, \emph{El Tepehua de Pisaflores}, was announced for the series Archivo de Lenguas Ind\'igenas de M\'exico. I did not confirm its publication: \url{https://www.ancientamericas.org/sites/default/files/01062MacKay01.pdf}
\item Carolyn MacKay's publications page (includes work on Pisaflores Tepehua phonology, IJAL 2013): \url{https://ballstate.academia.edu/CarolynMacKay}
\item Short basic Tepehua word list: \url{https://www.native-languages.org/tepehua_words.htm}
\item Endangered Languages Project entry: \url{https://www.endangeredlanguages.com/fr/node/40050}
\end{itemize}

\section{Religious and Community Resources}
\begin{itemize}
\item ScriptureEarth lists a Bible in Pisaflores Tepehua with text and audio: \url{https://www.scriptureearth.org/00eng.php?iso=tpp}
\end{itemize}
No community organisations were found in search results.

\section{Digital and Audio Resources}
Audio exists only as the scripture recordings above.

\section{Curriculum Development Context}
The documentation gives a base for orthography, vocabulary and grammar work. Any K-12 curriculum would need new materials developed with the speaker community. Check whether MacKay's lexicon and grammar are available before starting.
\end{document}